\documentclass[11pt,a4paper]{article}

\usepackage[T1]{fontenc}
\usepackage[utf8]{inputenc}
\usepackage[brazil]{babel}
\usepackage{lmodern}
\usepackage{microtype}
\usepackage[margin=2cm]{geometry}
\usepackage{hyperref}
\usepackage{titlesec}
\usepackage{xcolor}

\hypersetup{
  colorlinks=true,
  linkcolor=black,
  urlcolor=black,
  citecolor=black,
  pdfborder={0 0 0}
}

\definecolor{sectioncolor}{HTML}{111111}
\definecolor{accentcolor}{HTML}{1A1A1A}

\pagestyle{empty}
\setlength{\parindent}{0pt}
\setlength{\parskip}{5pt}
\linespread{1.05}
\selectfont

\titleformat{\section}
  {\bfseries\large\color{sectioncolor}}
  {}{0pt}{}[\vspace{-2pt}{\color{sectioncolor}\titlerule[0.8pt]}]
\titlespacing*{\section}{0pt}{12pt}{8pt}

\newcommand{\cvrole}[4]{%
  \vspace{4pt}%
  {\bfseries\color{accentcolor}#1}\hfill{\itshape #2}\\[-2pt]%
  {\itshape #3}\hfill #4\\[6pt]%
}

\begin{document}

\begin{center}
  {\fontsize{20}{24}\selectfont\bfseries Kamille Hartmann Maccagnan}\\[8pt]
  {\large PMO \textbar{} Governança de Projetos \textbar{} Indicadores e Stakeholders}\\[10pt]
  Campo Comprido, Curitiba-PR \enspace\textbullet\enspace (41) 9 9928-4424 \enspace\textbullet\enspace
  \href{mailto:kamillehartmannm@gmail.com}{kamillehartmannm@gmail.com}\\[3pt]
  \href{https://linkedin.com/in/kamille-hartmann-maccagnan/}{linkedin.com/in/kamille-hartmann-maccagnan/}
\end{center}

\vspace{6pt}

\section{Resumo Profissional}

Profissional com formação superior completa e pós-graduação em Liderança e Gestão em Tecnologia, com experiência em apoio a PMO, governança de projetos e acompanhamento de cronogramas, prazos, entregas e indicadores de desempenho em ambiente multinacional industrial. Atua na estruturação e atualização de informações de projetos, elaboração de relatórios executivos, monitoramento de KPIs e interface com stakeholders de áreas técnicas, operacionais e de negócio. Conhecimento em Boas Práticas PMOBOK, Scrum e Kanban, com domínio de Excel, PowerPoint, Power BI e ServiceNow. Perfil organizado, analítico e comunicativo, com visão sistêmica, proatividade e capacidade de atuar com múltiplas frentes de trabalho em ambientes dinâmicos.

\section{Experiência Profissional}

\cvrole{Anjuss}{Analista de Suporte de TI}{07/2025 -- Atual}

Atuo na organização e priorização de demandas, padronização de processos e documentação operacional, com criação e atualização de POPs que asseguram rastreabilidade, clareza e qualidade nas entregas. Lidero treinamentos de novos funcionários e apoio a rotinas administrativas e analíticas ligadas à operação. Desenvolvo relatórios e dashboards em Power BI para acompanhamento de indicadores e desempenho de processos, apoiando a identificação de desvios, pendências e oportunidades de melhoria. Participo da comunicação entre usuários e áreas técnicas, registrando ocorrências, encaminhamentos e alinhamentos necessários à execução das atividades.

\cvrole{HORSE do Brasil}{Estagiária de TI (IS/IT LATAM)}{07/2024 -- 07/2025}

Atuei no apoio a projetos e PMO em ambiente multinacional industrial, com acompanhamento de portfólio de projetos, cronogramas, prazos, entregas e responsáveis por iniciativas estratégicas. Atualizei o andamento de projetos por meio de relatórios executivos e dashboards em Power BI, apoiando reuniões de status, rituais de acompanhamento e a gestão de lideranças. Participei do levantamento de riscos, documentação de processos, análise de pendências e melhoria contínua, seguindo práticas de governança de projetos. Utilizei o ServiceNow para gestão de demandas e atuei como ponte entre áreas de negócio, operações, fornecedores e stakeholders internacionais, participando de reuniões em até três idiomas e registrando alinhamentos e encaminhamentos para apoio à tomada de decisão.

\cvrole{Restaurante Familiar}{Analista de Dados e Suporte Técnico}{01/2023 -- 07/2024}

Atuei em atividades analíticas e administrativas para apoio à gestão do negócio, com elaboração de relatórios, indicadores de desempenho e dashboards em Power BI para acompanhamento de resultados operacionais e financeiros. Organizei e padronizei informações para maior confiabilidade dos dados e apoio ao controle de performance. Atuei no alinhamento entre áreas operacionais e administrativas, contribuindo para o cumprimento de rotinas, prazos e melhoria da eficiência dos processos.

\section{Competências Técnicas}

\textbf{PMO e Gestão de Projetos:} Apoio a governança de projetos, acompanhamento de cronogramas e entregas, controle de prazos e escopo, portfólio de projetos, análise de riscos, reuniões de status, documentação de projetos, Boas Práticas PMOBOK, Scrum, Kanban, metodologias ágeis e híbridas.

\textbf{Indicadores e Relatórios:} KPIs, relatórios executivos e gerenciais, apresentações em PowerPoint, dashboards em Power BI, Excel avançado, análise de desvios e apoio à tomada de decisão.

\textbf{Stakeholders e Comunicação:} Interface com áreas técnicas, operacionais e de negócio, registro de decisões e encaminhamentos, trabalho em equipe multidisciplinar, comunicação com diferentes níveis hierárquicos, inglês avançado.

\textbf{Ferramentas:} ServiceNow, Trello, Monday, Power BI, Google Sheets, Pacote Office.

\textbf{Idiomas:} Inglês Avançado \enspace\textbullet\enspace Espanhol Básico \enspace\textbullet\enspace Coreano Básico

\section{Projetos e Conquistas}

1º lugar no Transformation Day Renault 2023 com solução orientada a melhoria de processos e gestão de entregas. Participação no Hackathon Bosch 2023 com foco em resolução de problemas de negócio. Acompanhamento de portfólio de projetos e elaboração de relatórios executivos em ambiente multinacional na HORSE. Estruturação de KPIs e dashboards para monitoramento de performance. Desenvolvimento de documentações e POPs com padronização e rastreabilidade.

\section{Formação Acadêmica}

\textbf{PUCPR} \hfill Curitiba, PR\\
Graduação em Gestão da Tecnologia da Informação \hfill Previsão de conclusão: 06/2026\\[6pt]

\textbf{Escola Conquer} \hfill Curitiba, PR\\
Pós-graduação em Liderança e Gestão em Tecnologia \hfill Conclusão: 2024\\[6pt]

\textbf{Universidade Tuiuti do Paraná} \hfill Curitiba, PR\\
Graduação em Medicina Veterinária \hfill 06/2019\\[6pt]

\textbf{Qualittas} \hfill Curitiba, PR\\
Pós-graduação em Clínica Médica e Cirúrgica de Animais Exóticos e Pets Não Convencionais \hfill 12/2022

\end{document}
